\section{Introduction}

\IEEEPARstart{T}{he} economics of consumer commerce in Southeast Asia are being reorganized around the \emph{delivery promise}. Shopee's in-house carrier SPX Express advertises same-day and instant fulfillment in Indonesian metros; quick-commerce operator Astro operates a dense network of dark stores that put a courier at the consumer's door in as little as 15--60 minutes. International precedents---Amazon's regionalized same-day and Freshippo store-warehouses, JD Logistics' integrated fulfillment estate, and Coupang's density-first Rocket Delivery---demonstrate a consistent endgame: \emph{whoever controls the fastest fulfillment loop accumulates the consumer relationship, the data, and ultimately the shelf}, including private-label brands that displace incumbent manufacturers.

ParagonCorp is Indonesia's largest beauty manufacturer, operating 16 brands (Wardah, Make Over, Emina, Kahf, Instaperfect, Crystallure, Labor\'e, Biodef, Tavi, Wonderly, Beyondly, OMG, Earth Love Life, Putri, and others) across roughly 2,000 active SKUs and bundles. Its current distribution stack---a National Distribution Centre (NDC) in Jakarta, a B2C DC in Tegal serving TikTok Shop and Shopee order flow, and a set of B2B DCs (Bogor, Bandung, Solo, Makassar, and others, including Malaysia) serving general- and modern-trade stores---was engineered for \emph{replenishment logistics}: moving cases to stores and parcels within days. It was not engineered for \emph{instant logistics}: moving a single eyeliner to a consumer in under an hour.

This asymmetry is now strategic. If marketplaces alone can meet the instant promise, ParagonCorp's brands become undifferentiated supply behind someone else's delivery experience. The competitive response specified here is not to out-build Shopee's parcel network, but to \emph{redefine the fulfillment layer for the beauty category}: a hyperlocal, brand-owned network of micro-fulfillment centers (MFCs) and consignment-store fulfillment nodes, fed by the existing NDC/DC backbone, replenished via digital mid-mile partners, and delivering through established on-demand fleets---so that \emph{any} channel order is fulfilled by the nearest ParagonCorp-controlled node, faster and cheaper per beauty basket than a generalist dark store can achieve.

\subsection{What Revision 5 Changes}

Revision~4 of this report was subjected to a structured six-lens design review before release. The review produced forty-eight candidate findings. Seven were discarded in an adversarial pass (Appendix~A-E). Of the forty-one that survived, thirteen required new or restated mathematics, fourteen required specification changes, and fourteen could not be resolved without facts nobody has confirmed and are registered as open items in Section~\ref{sec:open}, where they expand into twenty-five register rows because several name more than one quantity. Appendix~A gives the review method, the finding register, and the derivation of every correction with its supporting literature. The material changes are:

\begin{enumerate}
\item \textbf{The availability promise was arithmetically unreachable as specified.} Revision~4 set inventory targets from a per-SKU fill rate and promised 95\% attainment at the \emph{order} level. Because a PX1 order fails if any SKU in its basket is unavailable, the two are different quantities. Proposition~\ref{prop:orderfill} gives the inversion. Under an illustrative basket-size distribution with mean 2.45 distinct SKUs, a 98\% per-SKU target yields 95.2\% order availability before any queueing, picking, or travel variability is counted; at mean 3.05 SKUs it yields 94.0\%, below the promise.
\item \textbf{The service radius was computed from a quantile identity that does not hold.} Revision~4's Proposition~2 subtracted a single $\alpha$-quantile of in-node time from the promise budget and treated travel as deterministic at a mean speed. The PX1 constraint requires the $\alpha$-quantile of a \emph{sum}, and quantiles are not additive \cite{artzner1999,embrechts2002}. Proposition~\ref{prop:radius} gives a distribution-free replacement via an explicit risk-budget split. On illustrative parameters the feasible radius falls from 13.3\,km to 7.1\,km, raising the SLA node floor by a factor of $3.5$; whether that floor binds depends on parameters that are not yet measured.
\item \textbf{The risk-pooling penalty assumed a demand independence the design itself breaks.} Revision~4's Proposition~1 used the $\sqrt{n}$ safety-stock law, which holds only for independent zone demand. The design models live-commerce demand as spatially correlated Hawkes bursts and gives the distributed order manager cross-node substitution and rebalancing. Proposition~\ref{prop:pooling} restates the penalty under equicorrelated demand; Eppen's original result already covers this case \cite{eppen1979}. Theorem~\ref{thm:nodecount} is restated on the corrected law. On the illustrative parameters of Revision~4, moving from $\rho=0$ to $\rho=0.3$ raises the optimal node count from 9.1 to 13.2.
\item \textbf{Two theorems were cited outside their proven scope.} The $(1-1/e)$ coverage guarantee is proved for the uncapacitated case and was used to order a rollout in a network with hard node throughput caps; the real-time routing optimality result is proved for one wave and was used to justify continuous re-solving. Both scopes are now stated explicitly, and Section~\ref{sec:structural}-\ref{sec:capacitated} states what is and is not known for the capacitated variant.
\item \textbf{The oversell guarantee did not cover the operations the system performs.} Proposition~4 of Revision~4 assumed a single ledger partition; the distributed order manager rebalances stock \emph{across} partitions, and the store client may run offline for hours. Proposition~\ref{prop:atp} restates the invariant with a cross-partition commit rule and an effective-latency definition that accounts for the offline window.
\item \textbf{Staffing assumed Poisson arrivals into a wave-batched system.} Proposition~\ref{prop:staffing} restates square-root staffing under compound-Poisson (wave) arrivals. The inflation factor is $\sqrt{1+\lambda_f\Delta}$, which makes the wave length $\Delta$ an explicit trade between routing optimality and picker headcount.
\item \textbf{An acceptance gate tested model agreement rather than reality, two phases had no gate at all, and the program had no stop condition.} Section~\ref{sec:roadmap} re-cuts the Phase~1 gate against the design target, adds gates to Phases~3 and~4, adds a demand-side experiment (E7) and an incrementality holdout (E8), and states explicit kill criteria.
\item \textbf{Facts nobody has confirmed are no longer written as facts.} Twenty-five open items are tagged inline and registered in Section~\ref{sec:open}. Ten of them block Phase~0: eight must close before Phase~0 exits, and two before it starts.
\end{enumerate}

The contributions of this report are: (1) a state-of-the-art benchmark of global and regional instant-fulfillment architectures (Section~\ref{sec:bench}); (2) a formal mathematical definition of the ParaExpress network design problem, with theorems and propositions that justify the architecture and make it computable (Sections~\ref{sec:formal}--\ref{sec:structural}); (3) the ParaExpress reference network architecture and MFC operating model derived from those results (Sections~\ref{sec:arch}--\ref{sec:mfc}); (4) the business model and partnership topology (Section~\ref{sec:business}); (5) the solution specification (Section~\ref{sec:solution}); (6) the ParaTwin formal model, data contracts, and experiment plan (Section~\ref{sec:twin}); (7) a phased implementation roadmap with pilot acceptance criteria and kill criteria (Section~\ref{sec:roadmap}); and (8) an open-items register (Section~\ref{sec:open}) and a review-method appendix (Appendix~A).
